\documentclass[11pt,a4paper]{article}
\usepackage[utf8]{inputenc}
\usepackage{amsmath,amssymb,amsthm}
\usepackage{booktabs}
\usepackage{hyperref}
\usepackage{enumitem}
\usepackage[margin=1in]{geometry}

\newtheorem{theorem}{Theorem}
\newtheorem{lemma}[theorem]{Lemma}
\newtheorem{conjecture}[theorem]{Conjecture}
\newtheorem{proposition}[theorem]{Proposition}
\theoremstyle{definition}
\newtheorem{definition}[theorem]{Definition}
\theoremstyle{remark}
\newtheorem{remark}[theorem]{Remark}

\newcommand{\bOne}{\mathbf{1}}
\newcommand{\FF}{\mathbb{F}}
\newcommand{\ZZ}{\mathbb{Z}}
\newcommand{\NN}{\mathbb{N}}
\newcommand{\cS}{\mathcal{S}}
\newcommand{\cR}{\mathcal{R}}
\newcommand{\cC}{\mathcal{C}}
\newcommand{\cO}{\mathcal{O}}
\newcommand{\row}{\operatorname{row}}
\newcommand{\col}{\operatorname{col}}

\title{The Independent Set Inverse Problem\\for Bounded Treewidth Graphs}
\author{Igor Rivin \and with computational contributions from\\OpenAI o3, Google Gemini, and Anthropic Claude}
\date{March 2026}

\begin{document}
\maketitle

\begin{abstract}
We investigate which positive integers arise as the number of independent sets of a graph with treewidth at most~$k$. Using an algebraic framework that identifies series-parallel graph composition with matrix product and Hadamard product of $2\times 2$ signature matrices, we show computationally that every integer $n \geq 3$ is the independent set count of some series-parallel graph (treewidth $\leq 2$), verified for $n \leq 10{,}000$. We prove that the bisemigroup closure of the Fibonacci matrix modulo any prime $p$ is all of $M_2(\FF_p)$, eliminating all local obstructions. The surjectivity problem reduces, via a dot-product identity, to showing that a specific linear form on a structured subset of $\NN^2$ is cofinite.
\end{abstract}

\tableofcontents

%% ====================================================================
\section{Problem Statement}
%% ====================================================================

\textbf{Which positive integers arise as the number of independent sets of a graph with treewidth at most~$k$?}

This is a bounded-treewidth generalization of the classical IS inverse problem. For trees (treewidth~1), it is conjectured that all but finitely many positive integers are achievable, but no density result is known.

\subsection{Difficulty Hierarchy}

\begin{center}
\begin{tabular}{lrl}
\toprule
Treewidth & Achievable (of 10{,}000) & Mechanism \\
\midrule
$= 0$ (edgeless) & 13 (0.1\%) & Powers of 2 only \\
$\leq 1$ (forests) & 9{,}269 (92.7\%) & Tree DP + multiplicative closure \\
$\leq 2$ (series-parallel) & \textbf{10{,}000 (100\%)} & SP algebra (this paper) \\
$\leq 3$ & 10{,}000 (100\%) & Follows from tw $\leq 2$ \\
$\leq 4$ & All integers (proven) & Chain gadget / continued fraction \\
unbounded & All integers (trivial) & Threshold graphs / binary expansion \\
\bottomrule
\end{tabular}
\end{center}

\subsection{Main Results}

\begin{theorem}[Computational, verified for $n \leq 10{,}000$]
Every integer $n \geq 3$ is the number of independent sets of some series-parallel graph (treewidth $\leq 2$). Combined with the empty graph ($i = 1$) and a single vertex ($i = 2$), every positive integer is achievable at treewidth~$\leq 2$.
\end{theorem}

\begin{conjecture}
Every integer $n \geq 1$ is the number of independent sets of some graph with treewidth $\leq 2$.
\end{conjecture}

This makes the treewidth $\leq 3$ question (originally posed as the hard case between the open tree conjecture and the proven tw $\leq 4$ result) trivially resolved: treewidth~2 already suffices.

%% ====================================================================
\section{The Algebraic Framework}
%% ====================================================================

\subsection{Series-parallel graphs as matrix algebra}

A two-terminal series-parallel (SP) graph $G$ with terminals $s, t$ has an \emph{IS signature} $(f_{00}, f_{10}, f_{01}, f_{11})$, where $f_{ij}$ counts independent sets with $s$ in state $i$ and $t$ in state $j$. Arrange this as a $2 \times 2$ matrix:
\[
M_G = \begin{pmatrix} f_{00} & f_{01} \\ f_{10} & f_{11} \end{pmatrix}.
\]

The two SP composition operations become:
\begin{itemize}
  \item \textbf{Series composition} (identify inner terminals): $M_{G_1 ; G_2} = M_{G_1} \cdot M_{G_2}$ (matrix product).
  \item \textbf{Parallel composition} (identify both terminals): $M_{G_1 \| G_2} = M_{G_1} \circ M_{G_2}$ (Hadamard/entrywise product).
\end{itemize}

The IS count is the quadratic form:
\[
i(G) = \bOne^T M_G \bOne, \qquad \bOne = \begin{pmatrix}1\\1\end{pmatrix}.
\]

\subsection{The Fibonacci matrix}

The single edge $s$--$t$ has signature:
\[
M_E = \begin{pmatrix}1&1\\1&0\end{pmatrix}.
\]
This is the \emph{Fibonacci matrix}. As a M\"obius transformation, it acts as $z \mapsto 1 + 1/z$, with fixed point $\varphi = (1+\sqrt{5})/2$ (the golden ratio). It satisfies $M_E^2 = M_E + I$, has determinant $-1$, and lives in $GL_2(\ZZ)$.

Series powers give paths: $M_E^k$ has entries $(F_{k+1}, F_k; F_k, F_{k-1})$ and IS count $F_{k+2}$.

\subsection{The generated set}

Let $\cS$ be the smallest set of $2 \times 2$ nonneg integer matrices containing $M_E$ and closed under both matrix product and Hadamard product. Then:
\[
\{ i(G) : G \text{ is a series-parallel graph} \} \supseteq \{ \bOne^T M \bOne : M \in \cS \}.
\]

\textbf{Computational finding:} $\{ \bOne^T M \bOne : M \in \cS \} \supseteq \{3, 4, 5, \ldots, 10000\}$.

\subsection{The bisemigroup structure}

The set $\cS$ is closed under two associative binary operations with:
\begin{itemize}
  \item \textbf{Different identities}: $I$ for matrix product, $J = \left(\begin{smallmatrix}1&1\\1&1\end{smallmatrix}\right)$ for Hadamard product.
  \item \textbf{No distributive law}: neither operation distributes over the other.
  \item \textbf{Neither identity in $\cS$} (since $\cS$ is generated by $M_E$ alone).
\end{itemize}
This is a set with two independent semigroup structures (a \emph{bisemigroup}).

%% ====================================================================
\section{The Iterated Function System}
%% ====================================================================

\subsection{Linearization on $\ZZ^4$}

Identify each $2\times 2$ matrix with a vector $(a,b,c,d) \in \ZZ^4$ where $M = \left(\begin{smallmatrix}a&b\\c&d\end{smallmatrix}\right)$. The two operations (applied with $M_E$) become linear maps on $\ZZ^4$:

\textbf{Right-multiply by $M_E$ (series with edge):}
\[
R: (a,b,c,d) \mapsto (a+b,\; a,\; c+d,\; c).
\]
As a $4 \times 4$ matrix: $R = F \oplus F$, block-diagonal with the Fibonacci matrix acting independently on $(a,b)$ and $(c,d)$.

\textbf{Hadamard with $M_E$ (parallel with edge):}
\[
P: (a,b,c,d) \mapsto (a,\; b,\; c,\; 0).
\]
This is projection --- it simply kills the $(1,1)$ entry.

The IS count is the linear functional $\varphi(a,b,c,d) = a+b+c+d$.

\subsection{Properties of the generators}

\begin{center}
\begin{tabular}{lll}
\toprule
 & $R$ (series) & $P$ (parallel) \\
\midrule
Type & Expanding & Contracting (rank 3) \\
Eigenvalues & $\varphi, \varphi, -1/\varphi, -1/\varphi$ & $1, 1, 1, 0$ \\
Determinant & $1$ ($\in SL_4(\ZZ)$) & $0$ (not invertible) \\
\bottomrule
\end{tabular}
\end{center}

\subsection{Sequential collapse (Gemini)}

Both $R$ and $P$ are block-diagonal with the same $2 \times 2$ partition:
\[
R = F \oplus F, \qquad P = I_2 \oplus E, \qquad E = \begin{pmatrix}1&0\\0&0\end{pmatrix}.
\]
The top pair $(a,b)$ and bottom pair $(c,d)$ never mix. The bottom block collapses via $E F^k E = F_{k+1} \cdot E$. For a sequential word $X = R^{a_1} P R^{a_2} P \cdots P R^{a_{k+1}}$ (with $k$ projections and $N = \sum a_i$ total $R$'s), the IS count evaluates exactly to:
\[
f(X) = F_{N+3} + F_{a_1+2} \cdot F_{a_{k+1}+1} \cdot \prod_{i=2}^{k} F_{a_i+1}.
\]
Verified computationally for 9{,}330 words with 0 mismatches.

\textbf{Consequence:} The sequential IFS achieves only ${\sim}34\%$ of integers up to 1000. The sequential orbit has density tending to zero (like the primes), despite equidistribution mod~$p$ for all primes~$p$.

\begin{center}
\begin{tabular}{lr}
\toprule
Range & Sequential density \\
\midrule
$[3, 100]$ & 72\% \\
$[3, 500]$ & 46\% \\
$[3, 1000]$ & 34\% \\
$[3, 5000]$ & 19\% \\
$[3, 10000]$ & 13\% \\
\bottomrule
\end{tabular}
\end{center}

Positive density genuinely requires the branching Hadamard products.

%% ====================================================================
\section{The Local Theory}
%% ====================================================================

\subsection{The full closure theorem}

\begin{theorem}[Verified for $p = 2, 3, 5, 7, 11$]
\label{thm:full-closure}
The full bisemigroup closure of $M_E$ in $M_2(\FF_p)$ under matrix product and Hadamard product is:
\[
S_p = M_2(\FF_p).
\]
\end{theorem}

\begin{proof}[Proof sketch]
Define $K_u := \left(\begin{smallmatrix}u&1\\1&0\end{smallmatrix}\right)$. Then $K_1 = E$ and $K_2 = E^2 \circ E$ (Hadamard). Since $\det(E) = -1 \neq 0$, $E \in GL_2(\FF_p)$, and $E^{-1} \in S_p$ (as a power of $E$). The recurrence $K_{u+1} = K_2 \cdot E^{-1} \cdot K_u$ gives all $K_u$ for $u \in \FF_p$. From these:
\begin{align*}
U_t &:= K_{t+1} E^{-1} = \begin{pmatrix}1&t\\0&1\end{pmatrix}, \qquad
L_t := E^{-1} K_{t+1} = \begin{pmatrix}1&0\\t&1\end{pmatrix}.
\end{align*}
Diagonal matrices are obtained via Hadamard products with $I$ (which is a power of $E$):
\[
D_x := (U_1 L_{x-1}) \circ I = \begin{pmatrix}x&0\\0&1\end{pmatrix}, \qquad
D'_y := (L_{y-1} U_1) \circ I = \begin{pmatrix}1&0\\0&y\end{pmatrix}.
\]
The permutation matrix $P = E \circ E^{-1} = \left(\begin{smallmatrix}0&1\\1&0\end{smallmatrix}\right)$ is also in the closure. Any matrix with $a \neq 0$ factors as $L_{c/a} \, D_a \, D'_{d-bc/a} \, U_{b/a}$; remaining cases are handled by $P$ and zero diagonal factors.
\end{proof}

\begin{remark}[Sequential orbit is a proper subset]
The earlier claim that ``the orbit is a proper subset of $(\ZZ/m\ZZ)^4$'' was about the \emph{sequential} orbit (IFS applying $R$ or $P$ one at a time). The sequential orbit IS a proper subset (32.7\% of $M_2(\FF_7)$, decreasing with $p$). The full bisemigroup fills everything.
\end{remark}

\begin{remark}[Caveat for sieve methods]
The proof requires $E^{-1} = E^{\mathrm{ord}(E)-1}$, where $\mathrm{ord}(E)$ divides $|GL_2(\FF_p)| = p(p-1)^2(p+1)$. So the expression trees realizing arbitrary matrices mod~$p$ have depth polynomial in~$p$. This rules out the uniform-in-$p$ error bounds that sieve methods require. The archimedean problem (cofinite linear form on $\cR$) is genuinely independent of the local theory.
\end{remark}

\begin{center}
\begin{tabular}{lrr}
\toprule
$p$ & Sequential orbit & Full bisemigroup \\
\midrule
2 & 75\% of $M_2(\FF_2)$ & \textbf{100\%} \\
3 & 89\% & \textbf{100\%} \\
5 & 67\% & \textbf{100\%} \\
7 & 33\% & \textbf{100\%} \\
11 & --- & \textbf{100\%} \\
\bottomrule
\end{tabular}
\end{center}

%% ====================================================================
\section{The Global Theory: Dot-Product Reduction}
%% ====================================================================

\subsection{Transpose closure}

\begin{lemma}
\label{lem:transpose}
$\cS$ is closed under $M \mapsto M^T$, and consequently $\cR = \cC$.
\end{lemma}

\begin{proof}
$M_E$ is symmetric, $(A \circ B)^T = A^T \circ B^T$, and $(AB)^T = B^T A^T$. Any expression tree can be mirrored (swap children at each $\cdot$-node) to produce the transpose. Since $\col(M) = \row(M^T)$ and $M^T \in \cS$, we have $\cR = \cC$.
\end{proof}

\subsection{Dot-product identity}

\begin{lemma}
\label{lem:dot}
For $M, N \in \cS$,
\[
\varphi(MN) = \col(M) \cdot \row(N) = (M^T \bOne)^T (N \bOne),
\]
where $\row(N) = \bOne^T N \in \NN^2$ and $\col(M) = M^T \bOne \in \NN^2$.
\end{lemma}

\begin{proof}
$\varphi(MN) = \bOne^T MN \bOne = (M^T \bOne)^T (N \bOne)$.
\end{proof}

\begin{proposition}
$\varphi(\cS) \supseteq \{u \cdot v : u, v \in \cR\}$.
\end{proposition}

\subsection{The $(2, 1)$ shortcut}

Since $\col(M_E) = (2, 1)$, for any $N \in \cS$ with $\row(N) = (x, y)$:
\[
\varphi(M_E \cdot N) = 2x + y.
\]

\begin{conjecture}[Cofinite Linear Form]
\label{conj:clf}
The set $L := \{2x + y : (x, y) \in \cR\}$ is cofinite; equivalently, there exists $T$ such that $L \supseteq [T, \infty) \cap \ZZ$.
\end{conjecture}

Conjecture~\ref{conj:clf} plus a finite verification for $n < T$ would prove the main conjecture.

\subsection{The Fibonacci transform on row vectors}

Right-multiplication by $M_E$ acts on row vectors as the \emph{Fibonacci map}:
\[
\row(M \cdot M_E) = \row(M) \cdot M_E = (x + y,\; x).
\]
This map preserves $\gcd(x, y)$, and sends $2x + y \mapsto 3x + 2y = (2x+y) + x$. Repeated application gives $2x_n + y_n = F_{n+2}\, x + F_{n+1}\, y$, sweeping values that grow like~$\varphi^n$.

If the Hadamard product produces enough \emph{seed vectors} $(x, y)$ with modest~$y$, the Fibonacci transform propagates them into families where $2x + y$ sweeps intervals.

\subsection{Explicit seed families}

Two provable seed families:
\begin{enumerate}
  \item \textbf{Paths:} $\row(E^n) = (F_{n+2}, F_{n+1})$, giving $2x + y = 2F_{n+2} + F_{n+1}$.
  \item \textbf{Parallel paths:} $\row(E^m \circ E^n) = (F_{m+n+1}, F_{m+n-1})$, using the identity $F_{m+1}F_{n+1} + F_m F_n = F_{m+n+1}$. This gives $2x + y = 2F_{s+1} + F_{s-1}$ for $s = m + n$.
\end{enumerate}
Both families produce Fibonacci-spaced values, too sparse alone. The key question is whether additional families from deeper Hadamard products interleave to fill all gaps.

\subsection{Computational validation}

The dot-product identity was tested with the 4-iteration closure (2{,}594 signatures, 1{,}902 row/col vectors):

\begin{center}
\begin{tabular}{ll}
\toprule
Range & Dot product gaps \\
\midrule
$[1, 10{,}000]$ & Only $\{1, 2, 3, 4, 6, 9\}$ \\
$[10{,}001, 40{,}000]$ & None \\
$[40{,}001, 50{,}000]$ & 72 gaps (finite-depth artifact) \\
\bottomrule
\end{tabular}
\end{center}

Since $3, 4, 6, 9$ are all directly achievable as IS counts of small SP graphs, every integer in $[3, 10{,}000]$ is covered by the dot-product route alone.

For fixed col vector $(a, b)$, the achievable dot products become consecutive above a threshold:

\begin{center}
\begin{tabular}{lr}
\toprule
Col vector & Consecutive from \\
\midrule
$(4, 3)$ & 137 \\
$(3, 1)$ & 149 \\
$(3, 2)$ & 165 \\
$(2, 1)$ & 174 \\
\bottomrule
\end{tabular}
\end{center}

The row/col sum sets do \emph{not} form thick rectangles (coverage of $[1..10]^2$ is only 25\%), but their dot products are nearly surjective nonetheless.

%% ====================================================================
\section{Context: Why Standard Machinery Fails}
%% ====================================================================

\subsection{Arithmetic circuits and the branching bisemigroup}

Elements of $\cS$ are expression trees, not words. In an expression tree of depth $D$ with Hadamard nodes, entries are polynomials of degree up to $2^D$ in the Fibonacci numbers. This exponential degree explosion means:

\begin{center}
\begin{tabular}{lll}
\toprule
 & Sequential (words) & Bisemigroup (trees) \\
\midrule
Algebraic model & Random matrix product & Arithmetic circuit \\
Degree of entries & $O(1)$ (Fibonacci products) & $2^D$ (exponential) \\
Transfer operator & Linear, finite state mod $p$ & Nonlinear; no finite-state reduction \\
\bottomrule
\end{tabular}
\end{center}

\subsection{Barriers to equidistribution proofs}

\begin{enumerate}
  \item \textbf{Degree-characteristic barrier.} Weil/Deligne bounds require $d < p$; trees of depth $D \approx \log_2 p$ already exceed this.
  \item \textbf{PIT and circuit complexity.} Proving algebraic independence for these circuits is entangled with open problems in complexity theory.
  \item \textbf{No branching expander machine.} The measure evolution $\mu_{n+1} = \alpha(\mu_n * \mu_n) + \beta(\mu_n \circ \mu_n)$ is a nonlinear dynamical system with no precedent in the Bourgain--Gamburd framework.
\end{enumerate}

However, the full closure theorem (Theorem~\ref{thm:full-closure}) renders the local theory moot: $S_p = M_2(\FF_p)$ is established by an elementary algebraic argument, bypassing all of the above. The PIT/expander/sieve discussion is context for the general difficulty, not the proof strategy.

\subsection{The sieve gap}

Even equidistribution mod~$p$ with good rates is insufficient for surjectivity. Sieve methods need uniform-in-$p$ error bounds ($\sum_p \delta(p)/p < \infty$). The local theorem uses chains of length polynomial in~$p$, ruling out such uniformity. The archimedean problem is genuinely independent of the local theory.

\subsection{Analogy with Apollonian packings}

\begin{center}
\begin{tabular}{lll}
\toprule
 & Apollonian gasket & SP graph IS counts \\
\midrule
Space & Descartes quadruples $\in \ZZ^4$ & Matrix entries $\in \ZZ^4$ \\
Semigroup & Apollonian group (positivity) & $\langle R, P \rangle$ (non-invertible $P$) \\
Functional & Curvature & $\varphi = a+b+c+d$ \\
Local picture & Thin group in $O(3,1)$ & $S_p = M_2(\FF_p)$ (complete) \\
Positive density & Bourgain--Fuchs 2011 & Expected \\
Surjectivity & Open & Verified to 10{,}000 \\
\bottomrule
\end{tabular}
\end{center}

%% ====================================================================
\section{Prior Results}
%% ====================================================================

\subsection{The tw $\leq 4$ construction (proven universal)}

\textbf{(Due to OpenAI o3, March 2026.)} The ``chain gadget'' $C(m_1, \ldots, m_k)$ has IS count following the continued-fraction recurrence $A_i = m_i A_{i-1} + B_{i-1}$, $B_i = A_{i-1}$, with $i(G) = A_k + B_k$. The Euclidean algorithm guarantees parameters exist for any target~$T$, with $k = O(\log T)$. Each gadget has treewidth $\leq 2$; adding the spine gives treewidth $\leq 4$.

\subsection{The tree IS conjecture}

For trees (treewidth~1), it is conjectured that all but finitely many positive integers occur as $i(T)$. This is open. Trees miss ${\sim}7.3\%$ of integers up to 10{,}000 (including primes 7, 11, 19, 29, 31, \ldots). The jump from tw $\leq 1$ to tw $\leq 2$ closes all gaps; parallel composition (Hadamard product) has no tree analogue.

\subsection{Linek's theorem}

Linek (1989): every positive integer is the IS count of some bipartite graph, but with no treewidth bound.

%% ====================================================================
\section{Significance and Open Problems}
%% ====================================================================

\begin{enumerate}
  \item \textbf{Complete local picture.} The bisemigroup closure mod~$p$ is all of $M_2(\FF_p)$, eliminating all local obstructions. This is far stronger than $\varphi$-equidistribution.

  \item \textbf{Clean reduction to 1D number theory.} Via transpose closure ($\cR = \cC$) and the dot-product identity, the surjectivity problem reduces to: is $\{2x + y : (x, y) \in \cR\}$ cofinite?

  \item \textbf{Finite boundary algebra viewpoint.} The $2 \times 2$ terminal signatures are the $k = 2$ instance of the graph-homomorphism / Holant framework. The Holant clone theory and Fibonacci-gate literature~\cite{arXiv:1811.00817} may provide structural lemmas.

  \item \textbf{The remaining open problem} is purely archimedean: prove Conjecture~\ref{conj:clf} (the cofinite linear form conjecture for the row-sum set $\cR$). This is a statement in elementary additive number theory about a structured subset of $\NN^2$ generated by the Fibonacci map and entrywise products.
\end{enumerate}

%% ====================================================================
\begin{thebibliography}{99}

\bibitem{BGS}
J.~Bourgain, A.~Gamburd, P.~Sarnak,
\emph{Affine linear sieve, expanders, and sum-product},
Invent.\ Math.\ \textbf{179} (2010), 559--644.

\bibitem{BF}
J.~Bourgain, E.~Fuchs,
\emph{A proof of the positive density conjecture for integer Apollonian circle packings},
J.~Amer.\ Math.\ Soc.\ \textbf{24} (2011), 945--967.

\bibitem{BK}
J.~Bourgain, A.~Kontorovich,
\emph{On the local-global conjecture for integral Apollonian gaskets},
Invent.\ Math.\ \textbf{196} (2014), 589--650.

\bibitem{Linek}
V.~Linek,
\emph{Bipartite graphs can have any number of independent sets},
Discrete Math.\ \textbf{76} (1989), 131--136.

\bibitem{LM}
V.~E.~Levit, E.~Mandrescu,
\emph{The independence polynomial of a graph --- a survey},
Proceedings of the 1st International Conference on Algebraic Informatics, 2005.

\bibitem{SGV}
A.~Salehi Golsefidy, P.~P.~Varj\'u,
\emph{Expansion in perfect groups},
Geom.\ Funct.\ Anal.\ \textbf{22} (2012), 1832--1891.

\bibitem{arXiv:1811.00817}
J.-Y.~Cai, Z.~Fu, H.~Guo, T.~Williams,
\emph{A Holant dichotomy: is the FKT algorithm universal?},
arXiv:1811.00817, 2018.

\end{thebibliography}

\end{document}
